\begin{theorem}[Cauchy–Schwarz Inequality]
For real numbers $a_{1},\dots ,a_{n}$ and $b_{1},\dots ,b_{n}$,
\[
\Bigl(\sum_{i=1}^{n} a_{i}b_{i}\Bigr)^{2}
\le
\Bigl(\sum_{i=1}^{n} a_{i}^{2}\Bigr)\Bigl(\sum_{i=1}^{n} b_{i}^{2}\Bigr),
\]
with equality iff there exists $\lambda\in\mathbb{R}$ such that $a_{i}=\lambda b_{i}$ for every $i$.
\end{theorem}
\begin{proof}
\begin{enumerate}
\item \textbf{Interpretation as vectors and inner product.}
Let
\[
a=(a_{1},\dots ,a_{n}),\qquad b=(b_{1},\dots ,b_{n})\in\mathbb{R}^{n}.
\]
On $\mathbb{R}^{n}$ consider the standard inner product
\[
\langle a,b\rangle:=\sum_{i=1}^{n} a_{i}b_{i},
\]
and the induced Euclidean norm
\[
\|a\|:=\sqrt{\langle a,a\rangle}= \sqrt{\sum_{i=1}^{n} a_{i}^{2}},\qquad
\|b\|:=\sqrt{\sum_{i=1}^{n} b_{i}^{2}}.
\]

\item \textbf{Non‑negativity of a squared norm.}
For any real scalar $t$ the vector $a-tb$ satisfies $\|a-tb\|^{2}\ge0$.
Expanding using the inner product gives
\begin{align*}
\|a-tb\|^{2}
&=\langle a-tb,\,a-tb\rangle \\
&=\langle a,a\rangle-2t\langle a,b\rangle+t^{2}\langle b,b\rangle \\
&=\sum_{i=1}^{n} a_{i}^{2}-2t\sum_{i=1}^{n} a_{i}b_{i}
   +t^{2}\sum_{i=1}^{n} b_{i}^{2}.
\end{align*}

\item \textbf{Quadratic in $t$.}
Define
\[
Q(t):=\Bigl(\sum_{i=1}^{n} b_{i}^{2}\Bigr)t^{2}
       -2\Bigl(\sum_{i=1}^{n} a_{i}b_{i}\Bigr)t
       +\Bigl(\sum_{i=1}^{n} a_{i}^{2}\Bigr).
\]
From the previous step $Q(t)=\|a-tb\|^{2}\ge0$ for every real $t$.
If $\sum b_{i}^{2}=0$ then $b=0$ and the inequality is trivial, so we may assume the leading coefficient is positive.

\item \textbf{Discriminant condition.}
A quadratic $\alpha t^{2}+\beta t+\gamma$ with $\alpha>0$ that is non‑negative for all $t\in\mathbb{R}$ must satisfy $\beta^{2}-4\alpha\gamma\le0$.
Here
\[
\alpha=\sum b_{i}^{2},\qquad
\beta=-2\sum a_{i}b_{i},\qquad
\gamma=\sum a_{i}^{2}.
\]
Hence
\begin{align*}
\Delta
&=\beta^{2}-4\alpha\gamma \\
&=\bigl(-2\sum a_{i}b_{i}\bigr)^{2}
   -4\Bigl(\sum b_{i}^{2}\Bigr)\Bigl(\sum a_{i}^{2}\Bigr) \\
&=4\Bigl[\bigl(\sum a_{i}b_{i}\bigr)^{2}
        -\bigl(\sum a_{i}^{2}\bigr)\bigl(\sum b_{i}^{2}\bigr)\Bigr].
\end{align*}
Since $\Delta\le0$, we obtain
\[
\bigl(\sum a_{i}b_{i}\bigr)^{2}
\le
\bigl(\sum a_{i}^{2}\bigr)\bigl(\sum b_{i}^{2}\bigr),
\]
which is precisely the Cauchy–Schwarz inequality.

\item \textbf{Equality case.}
Equality holds iff $\Delta=0$, i.e. iff the minimizing
\[
t_{0}= \frac{\sum a_{i}b_{i}}{\sum b_{i}^{2}}
\]
satisfies $\|a-t_{0}b\|^{2}=0$. This is equivalent to $a=t_{0}b$, or $a_{i}=t_{0}b_{i}$ for every $i$.
Conversely, if $a_{i}=\lambda b_{i}$ for some $\lambda\in\mathbb{R}$, then both sides of the inequality coincide. Hence equality occurs exactly when the vectors $a$ and $b$ are linearly dependent.
\end{enumerate}
\end{proof}